\section{Implementation, evidence format, and conservative dispatch}
\label{sec:implementation}

\subsection{The implemented operation}

The new producer is \code{fastunknot/singleton_dag.py}. Its planner stores
one immutable pair per grammar node: the maximum absolute generator label
and that label's occurrence count, capped at two. Concatenation retains the
larger maximum, or adds the capped counts when the maxima agree. Each
current root is inspected in its original slot; the first eligible donor
for each child is chosen. The output is sorted by child label, making its
dependency order explicit.

The cache contains word content only. Relator slots and current live-label
membership are checked afresh. This matters because roots can move between
slots, normalization creates different roots, and a previously live
generator can have been eliminated. Planning neither mutates the
presentation nor allocates substituted images.

Application builds the prefix and suffix surrounding each unique pivot.
The inverse required by a positive pivot is represented by a virtual sign.
A single iterative traversal then evaluates signed grammar states, with
selected terminals redirecting to their definitions. An inverse
concatenation reverses the order of its children. Each signed state is
evaluated at most once. Both images of a definition are shared when needed;
they are not independently expanded or repeatedly inverted.

The helper also supports arbitrary internally supplied acyclic schedules.
An ordered schedule uses rank maxima and capped counts, regardless of its
numeric labels. An out-of-order internal schedule can use a generic
presence/repetition mask fallback and a cycle check. That fallback has a
larger polynomial metadata cost. The public certificate path always uses
an ordered witness and needs no dense support masks.

\subsection{Version-eight certificates}

A batch move contains only its kind and an ordered pivot list:
\begin{lstlisting}
{
  "kind": "singleton_dag",
  "pivots": [
    {"relation": 0, "generator": 2},
    {"relation": 1, "generator": 3}
  ]
}
\end{lstlisting}
The order is part of the evidence. All selected generators in a donor's
right-hand side must occur earlier in this list. Surviving generators have
rank zero. The certificate gives no replacement word, inferred exponent,
dependency list, source grammar node identifier, or claimed hash equality
for the checker to trust.

The complete group certificate still binds the input planar diagram,
method, status, moves, and terminal evidence. A trace using this move has
version eight. Both source replayers reject it if relabelled as an older
format. Previously supported versions remain accepted through their
existing interpretation.

\subsection{Independent compressed and literal replay}

The new checker module, \code{singleton_dag_verify.py}, imports no producer
helper. It requires exact field sets, actual integers rather than boolean
substitutes, valid original slots, live pivot labels, distinct slots, and
distinct pivots. It independently assigns certificate ranks and verifies
that each donor has its own pivot as a unique maximum-rank occurrence.
This checks singleton status and dependency order in one pass.

Its compressed implementation locates the pivot position, obtains the two
source substrings, and builds its own signed substitution evaluation.
Its context construction differs from the producer's direct sibling
accumulation. The literal implementation instead extracts explicit words,
builds and checks the dependencies, and substitutes definitions in
topological order with allocation checks before growth. The latter is an
independent oracle for bounded-size examples, not a polynomial method in
compressed input size when the expanded output is enormous.

Every definition is checked even when all its occurrences would lie in
discarded donor slots. This prevents an unused cyclic system from escaping
validation. Both replayers retain all non-donor slots. Rejection leaves
the root array and live-generator set unchanged; a compressed arena may
have allocated temporary nodes before rejection or resource exhaustion.
No claim of rollback of internal cache allocation is needed for soundness.

The outer checker independently reconstructs all source Wirtinger relations
before reading the moves. It finishes a rank-one trace only after checking
every retained root's exponent. Cyclicity is not inferred from one selected
relation or from the number of remaining labels alone.

\subsection{Why the delivered dispatcher is terminal-only}

The first integration applied every discovered batch of at least two
pivots. That policy was algebraically sound and worked on many actual
diagrams. It also changed the residual spelling on the Gordian fixture
enough to make the subsequent heuristic search hit a large work cap.
The local node bound does not control the difficulty of later Whitehead
or relator operations on that residual state.

The delivered public option therefore applies a batch only if
\begin{equation}
 2\leq |C|=|X|-1.
 \label{eq:terminal-guard}
\end{equation}
If the guard fails, planning has not changed the relators, their slots,
the live labels, or the established move order of the same compressed
backend with the same other flags. Search continues from the
intact state. If it succeeds, the batch leads directly to the existing
exact rank-one exponent check. The general batch operation and its proof
remain useful to research clients; the dispatcher makes a narrower choice
about when to use it during heuristic search.

This guard has a simple logical property. For a source knot presentation
and an uncapped correct preceding trace, an exact elimination to rank one
must have zero exponent in every retained relation: otherwise the resulting
one-generator presentation would have finite abelianization, contradicting
the knot group's abelianization $\Z$. The checker still verifies the
exponents explicitly. The public option can thus finish at the guarded
state without requiring any subsequent search heuristic to be effective.

Planning consumes resources even when it skips a batch. Therefore the
unchanged-state property does not imply identical completion under every
finite work or time allowance. The feature stays opt-in, and the experiments
record capped nondecisions and overhead. No universal wall-time dominance
over the predecessor is claimed.

\subsection{Public entry points and compatibility}

The group-stage interfaces accept \code{singleton_dag=True}; this selects
compressed search and independent compressed replay. The whole recognizer
accepts \code{group_singleton_dag=True} together with \code{use_group=True}.
The command line exposes \code{--group-singleton-dag}, which enables the
group stage automatically. For example:
\begin{lstlisting}[language=Python]
from fastunknot import Diagram
from fastunknot.group_certificate import group_decide

d = Diagram.from_braid(129, list(range(1, 129)))
result = group_decide(d, singleton_dag=True, seconds=None)
assert result["status"] == "UNKNOT"
assert result["certificate"]["version"] == 8
\end{lstlisting}

The option defaults to false. Existing projection and forest flags retain
their original meaning. Source replay accepts old certificates, and an audit
compares all three old modes against the entire pinned package. The Rust
port is not changed in this package. The new Python implementation uses the
standard library and requires the same Python language baseline as the
maintained package.

\subsection{Resource model}

The same work, node, initial-letter, and caller-cancellation mechanisms as
the surrounding compressed search remain in force. The new operation
does not use expanded word length as a loop bound. Its counters separately
report planner attempts, new immutable descriptors, candidate slots,
unique-occurrence path steps, signed evaluation states, batch pivots,
newly allocated nodes, and skipped partial batches.

The existing \code{projection_normalizations} counter counts raw-to-normal
handoffs in the common search loop. A terminal singleton batch causes none.
In the retained eager prototype, handoffs after a partial singleton batch
use that existing counter name; the result files preserve the original
names rather than silently recategorizing them.

Charged work is an implementation accounting measure. It includes explicit
charges such as the pivot sort and is useful for diagnosing scaling, but
it is neither a hardware instruction count nor a substitute for the
separate bit-complexity proof and completed-call timings.
